\subsection{The rank-one stationary family near the exceptional endpoint}
\label{sec:rank-one-stationary-family}

We now construct an analytic family of rank-one bipodal
candidates.  The numerical exploration in \Cref{fig:bipodal-parameter-maps}
suggests that rank-one graphons describe the optimizers along a curve of
parameter pairs $(p,r)$ approaching $P_*$.  Motivated by this observation,
we seek rank-one solutions of the full graphon stationarity equations,
together with the block-proportion stationarity condition.
At this stage, the resulting family consists only of stationary
rank-one bipodal candidates; the subsequent arguments establish their optimality and
uniqueness along the corresponding parameter curve.

Throughout this subsection, graphons are assumed to take values in a compact
subinterval of $(0,1)$.  For a rank-one graphon $W=f\otimes f$ with
measurable $f:[0,1]\to[0,1]$, this requires that, for some $\eta>0$,
\[
  \eta\le f\le1-\eta
  \qquad\text{a.e.}
\]
This interiority condition is used only to construct the candidate family in
this subsection.  It ensures that $W+\varepsilon U$ remains a graphon for every
bounded symmetric function $U$ and all sufficiently small $|\varepsilon|$, so
that stationarity can be tested in every bounded direction.  The subsequent
optimality and uniqueness arguments impose no such restriction on competing
graphons.  We say that $W$ is \emph{stationary} for
$I_p-\mu t(H,\,\cdot\,)$ if
\begin{equation}\label{eq:graphon-stationarity}
  \left.\frac{d}{d\varepsilon}\right|_{\varepsilon=0}
  \left[
    I_p(W+\varepsilon U)
    -\mu t(H,W+\varepsilon U)
  \right]=0
\end{equation}
for every such $U$.
For a rank-one graphon $W=f\otimes f$, we say that $W$ \emph{satisfies the
KKT conditions} at $(p,r)$ with multiplier $\mu$, if \eqref{eq:graphon-stationarity}
holds with $\mu\ge0$ and $t(H,W)=r^m$.  This is the graphon
analog of the usual first-order Karush--Kuhn--Tucker condition for
constrained finite-dimensional optimization.

To formulate the separate first-order condition for block proportion,
for $\alpha\in(0,1)$ and $s,t\in(0,1)$, set
\[
  f_{\alpha,s,t}
  :=s\1_{[0,\alpha]}+t\1_{(\alpha,1]},
  \qquad
  W_{\alpha,s,t}:=f_{\alpha,s,t}\otimes f_{\alpha,s,t}.
\]
Changing $\alpha$ moves the boundary between the two vertex classes.  Unlike
the perturbations $W_{\alpha,s,t}+\varepsilon U$ used in graphon stationarity,
replacing $\alpha$ by $\alpha+\varepsilon$, when $s\ne t$, changes the graphon
by an amount bounded away from zero on a set whose measure tends to zero.
Hence this variation is not small in $L^\infty$ and is not covered by graphon
stationarity.  We therefore impose the separate condition
\begin{equation}\label{eq:block-stationarity}
  \left.\frac{d}{d\beta}\right|_{\beta=\alpha}
  \left[
    I_p(W_{\beta,s,t})-\mu t(H,W_{\beta,s,t})
  \right]=0,
\end{equation}
obtained by varying the block proportion while keeping $s$ and $t$ fixed.

The following lemma makes a structural reduction: it shows that bipodality is
forced by the KKT conditions within the interior rank-one class.  The
construction lemma later begins by seeking a bipodal family; this structural
result justifies that choice and removes the need to assume bipodality in
the accompanying local uniqueness assertion.  The proof reduces the stationarity condition
\eqref{eq:graphon-stationarity} to a scalar equation and bounds the
number of its roots.  For this reduction, define the \emph{$d$-th moment} of
$f$ by
\[
  q(f):=\int_0^1 f(x)^d\,dx.
\]
Note that any rank-one graphon $W=f\otimes f$ satisfies
$t(H,W)=q(f)^v$ depending only on $H$ through $d$ and $v$:
\[
  t(H,f\otimes f)
  =\int_{[0,1]^v}\prod_{(i,j)\in E(H)}f(x_i)f(x_j)\,dx_1\cdots dx_v
  =\int_{[0,1]^v}\prod_{i=1}^v f(x_i)^{d}\,dx_1\cdots dx_v
  =q(f)^v.
\]

\begin{lemma}[Bipodality of stationary rank-one graphons]
\label{lem:stationary-rank-one-bipodality}
Let $p,r\in(0,1)$, let $\mu\ge0$, and let
$f:[0,1]\to[0,1]$ be measurable.  Suppose that, for some $\eta>0$,
\[
  \eta\le f\le1-\eta
  \qquad\text{a.e.},
\]
and that $f\otimes f$ satisfies the KKT conditions at $(p,r)$ with multiplier
$\mu$.  If $f$ is not a.e.\ constant, then there exist $\alpha\in(0,1)$ and
$0<s<t<1$ such that, after a measure-preserving relabeling,
\[
  f=f_{\alpha,s,t}
  \qquad\text{a.e.}
\]
\end{lemma}

\begin{proof}
Set $\gamma:=\mu m q(f)^{v-2}$ and $F_{p,\gamma}(z):=J_p'(z)-\gamma z^{d-1}$.
We first claim that
\begin{equation}
\label{eq:rank-one-kkt}
  F_{p,\gamma}(f(x)f(y))=0
  \quad\text{for a.e. }(x,y)\in[0,1]^2.
\end{equation}
Let $U$ be any bounded symmetric function on $[0,1]^2$ and set
$W_\varepsilon=f\otimes f+\varepsilon U$.  The variation of $I_p$ is
\[
  \left.\frac{d}{d\varepsilon}I_p(W_\varepsilon)\right|_{\varepsilon=0}
  =\iint J_p'(f(x)f(y))U(x,y)\,dx\,dy.
\]
Since $H$ is $d$-regular, differentiating one of the $m$ edge factors in
$t(H,W_\varepsilon)$ gives
\[
  \left.\frac{d}{d\varepsilon}t(H,W_\varepsilon)\right|_{\varepsilon=0}
  =\iint m q(f)^{v-2}f(x)^{d-1}f(y)^{d-1}
  U(x,y)\,dx\,dy.
\]
Stationarity for every $U$ therefore yields
\[
  J_p'(f(x)f(y))
  =\mu m q(f)^{v-2}f(x)^{d-1}f(y)^{d-1}
  =\gamma\bigl(f(x)f(y)\bigr)^{d-1}
\]
for almost every $(x,y)$, proving
\eqref{eq:rank-one-kkt}.

We next show that $F_{p,\gamma}(z)=0$ has at most three
solutions for $z\in(0,1)$.  Since $J_p''(z)=1/(z(1-z))$,
\[
  F_{p,\gamma}'(z)
  =\frac1{z(1-z)}-\gamma(d-1)z^{d-2}.
\]
If $\gamma\le0$, then $F_{p,\gamma}'>0$ on $(0,1)$, so
$F_{p,\gamma}$ has at most one zero.  If $\gamma>0$, the critical-point
equation is
\[
  z^{d-1}(1-z)=\frac1{(d-1)\gamma}.
\]
The function $z\mapsto z^{d-1}(1-z)$ is strictly increasing on
$(0,(d-1)/d)$ and strictly decreasing on $((d-1)/d,1)$.  Hence
$F_{p,\gamma}'$ has at most two distinct zeros,
and by Rolle's theorem, $F_{p,\gamma}$ has at most three distinct zeros.

By Fubini's theorem, one can choose $y_0\in[0,1]$ such that
\[
  F_{p,\gamma}(f(x)f(y_0))=0
  \qquad\text{for a.e. }x.
\]
The lower bound on $f$ gives $f(y_0)>0$.  Hence multiplication by
$f(y_0)$ is injective, and the three-root bound for $F_{p,\gamma}$ implies
that $f$ takes at most three values up to null sets.

Suppose that it takes three distinct values $0<a<b<c<1$ up to null sets.
Since \eqref{eq:rank-one-kkt} holds outside a null subset
of $[0,1]^2$, each of
\[
  a^2,\qquad ab,\qquad ac,\qquad bc,\qquad c^2
\]
is a zero of $F_{p,\gamma}$.  They satisfy
$a^2<ab<ac<bc<c^2$, contradicting the three-root bound.  Hence $f$ takes at
most two values.
Therefore, if $f$ is not a.e. constant, it has the form
\[
  f=s\1_A+t\1_{A^c}
\]
for some $0<s<t<1$ and measurable $A$ with $0<\lambda(A)<1$.  Setting
$\alpha:=\lambda(A)$ and relabeling $A$ as $[0,\alpha]$ gives
$f=f_{\alpha,s,t}$ almost everywhere.
\end{proof}

\Cref{lem:stationary-rank-one-bipodality} reduces the problem to a finite-dimensional
one: after a measure-preserving relabeling, the factor of every nonconstant
rank-one graphon satisfying the KKT conditions has the form
$f_{\alpha,s,t}$ for some $\alpha\in(0,1)$ and $0<s<t<1$.  It does not,
however, construct such graphons near $P_*$ or determine their parameters.
The next lemma constructs this analytic family.
Reparametrize $s$ and $t$ by
\[
  u:=\frac{s+t}{2},
  \qquad
  h:=\frac{t-s}{2},
  \qquad
  \gamma:=\mu m q(f_{\alpha,s,t})^{v-2}.
\]
Thus, $u$ is the midpoint of the two factor values, $|h|$ is half their
separation, and $\gamma$ is the coefficient appearing in the scalar
stationarity equation \eqref{eq:rank-one-kkt}.
The three stationarity equations, evaluated at the three graphon values,
determine $u,p,\gamma$; the block-proportion equation then determines
$\alpha$, and the defining relation for $\gamma$ determines the multiplier $\mu$.
Changing $h$ to $-h$ exchanges the two factor values and their proportions,
while remaining quantities are unchanged.

We denote the value of the coefficient $\gamma$ at the singular endpoint by
\[
  \gamma_*:=\left(\frac{d}{d-1}\right)^d=r_*^{-d}.
\]

\begin{lemma}[Analytic rank-one KKT family]
\label{lem:rank-one-kkt-family}
There exist $h_0>0$ and real-analytic functions
\[
  h\longmapsto(u_h,p_h,\gamma_h,\alpha_h)
  \qquad (|h|<h_0)
\]
such that
\[
  (u_0,p_0,\gamma_0,\alpha_0)
  =\left(u_*,p_*,\gamma_*,\frac12\right),
\]
and the following properties hold.  First, for all $|h|<h_0$,
\[
  u_{-h}=u_h,\qquad p_{-h}=p_h,\qquad
  \gamma_{-h}=\gamma_h,\qquad \alpha_{-h}=1-\alpha_h.
\]
Next, for $|h|<h_0$, define the two-valued factor and its rank-one graphon by
\[
  s_h:=u_h-h,
  \qquad 
  t_h:=u_h+h,
  \qquad
  f_h:=f_{\alpha_h,s_h,t_h}, 
  \qquad
  W_h:=f_h\otimes f_h,
\]
and define the $d$-th moment of $f_h$, the associated target density, and the
Lagrange multiplier by
\[
  q_h:=q(f_h)=\alpha_hs_h^d+(1-\alpha_h)t_h^d,
  \qquad
  r_h:=q_h^{2/d}, 
  \qquad
  \mu_h:=\frac{\gamma_h}{m q_h^{v-2}}.
\]
The resulting functions $q_h,r_h,\mu_h$ are real analytic, and for all
$|h|<h_0$, these quantities satisfy
\[
  0<p_h<r_h<1,
  \qquad s_h,t_h\in(0,1),
  \qquad \alpha_h\in(0,1),
  \qquad \mu_h>0,
  \qquad\text{and}\qquad
  t(H,W_h)=r_h^m.
\]
Finally, for $0<|h|<h_0$, the nonconstant rank-one graphon $W_h$ satisfies
the KKT conditions at $(p_h,r_h)$ with multiplier $\mu_h$, and
\eqref{eq:block-stationarity} holds.  Conversely, in the precise
sense stated below, let $W=f\otimes f$ be a nearby nonconstant rank-one
graphon whose factor $f$ is bounded away from $0$ and $1$.  If $W$ satisfies
the KKT conditions with multiplier $\mu$ and
\eqref{eq:block-stationarity} holds, then $W$ agrees, up to
measure-preserving relabeling, with a unique member $W_h$ indexed by
$h\in(0,h_0)$, and $\mu=\mu_h$.
\end{lemma}

The graphons $W_h$ supplied by \Cref{lem:rank-one-kkt-family},
with their measure-preserving relabelings, form the family used in
\Cref{thm:singular-endpoint}.  Before proving the lemma, we clarify the
family at the singular endpoint and the local uniqueness assertion.
At $h=0$, the definitions in the lemma give
\[
  s_0=t_0=u_*,
  \qquad
  q_0=u_*^d,
  \qquad
  r_0=u_*^2=r_*,
  \qquad
  \mu_0=\frac1{m r_*^m},
  \qquad
  W_0\equiv r_*.
\]
Thus, $(p_0,r_0)=P_*$, so the family passes through the constant graphon
$W_0\equiv r_*$ at the singular endpoint.  Also, the two factor values $s_0$ and
$t_0$ coincide, as do the three values $s_0^2,s_0t_0,t_0^2$ taken by the
graphon $W_0$.
For the local uniqueness assertion, consider a graphon $W=f\otimes f$ and a
multiplier $\mu$ appearing in the converse assertion of the lemma.  By
\Cref{lem:stationary-rank-one-bipodality}, after a measure-preserving relabeling its
factor has the form $f_{\alpha,s,t}$ with
$\alpha\in(0,1)$ and $0<s<t<1$.  If
$(s,t,p,r,\mu,\alpha)$ is sufficiently close to
\[
  \left(u_*,u_*,p_*,r_*,\frac1{m r_*^m},\frac12\right),
\]
then $h:=(t-s)/2$ belongs to $(0,h_0)$ and
\[
  (s,t,p,r,\mu,\alpha)
  =(s_h,t_h,p_h,r_h,\mu_h,\alpha_h).
\]

\begin{proof}[Proof outline]
Write
\[
  \ell(z):=\log\frac{1-z}{z},
  \qquad
  \ell_*:=\ell(p_*)=\frac{d}{d-1}-\log(d-1).
\]
By \Cref{lem:stationary-rank-one-bipodality}, every nonconstant rank-one graphon
satisfying the KKT conditions has, after a measure-preserving relabeling, a
two-valued factor.  Motivated by this reduction, we seek a family of the form
\[
  s:=u-h,
  \qquad
  t:=u+h,
  \qquad
  f:=f_{\alpha,s,t},
  \qquad
  W:=f\otimes f.
\]
Since $f\otimes f$ takes the three values $s^2,st,t^2$, the scalar
stationarity equation \eqref{eq:rank-one-kkt} requires
\begin{equation}\label{eq:three-value-kkt}
  \ell(p)-\ell(s^2)=\gamma s^{2d-2},\qquad
  \ell(p)-\ell(st)=\gamma(st)^{d-1},\qquad
  \ell(p)-\ell(t^2)=\gamma t^{2d-2}.
\end{equation}
We therefore begin by solving this necessary system for $u,p,\gamma$ as
functions of $h$; the block proportion $\alpha$ does not occur and will be
determined later.
At $h=0$, these equations have the solution $(u,p,\gamma)=(u_*,p_*,\gamma_*)$.
For $h\ne0$, subtracting consecutive equations in \eqref{eq:three-value-kkt} eliminates
$\ell(p)$ and expresses $\gamma$ as each of the two divided slopes across
$s^2,st,t^2$.  Although each slope appears as $0/0$ at $h=0$, both extend
analytically.  Let $A(u,h)$ denote the difference of the two divided slopes,
so the system requires $A(u,h)=0$.  The symmetry $h\mapsto-h$ makes $A$ odd in $h$,
and hence $A(u,h)=hB(u,h)$ for some analytic $B$.  A direct calculation shows that
$B(u,0)$ has the unique nondegenerate zero $u=u_*$.  The analytic implicit
function theorem therefore gives $h_0>0$ and a unique even analytic solution
$u=u_h$ for $|h|<h_0$ satisfying $B(u_h,h)=A(u_h,h)=0$.  The common divided slope then defines the even function
$\gamma=\gamma_h$, and the middle equation in
\eqref{eq:three-value-kkt} determines the even function $p=p_h$.

We now choose the block proportion $\alpha$ as a function of $h$.
Once $\alpha$ is fixed, set
\[
  q:=\alpha s_h^d+(1-\alpha)t_h^d,
  \qquad
  r:=q^{2/d},
  \qquad
  \mu:=\frac{\gamma_h}{m q^{v-2}},
  \qquad
  f:=f_{\alpha,s_h,t_h},
  \qquad
  W:=f\otimes f.
\]
Also define the function
\[
  \mathcal M_h(z):=J_{p_h}(z)-\frac{\gamma_h}{d}z^d
  \qquad (0 < z < 1).
\]
Evaluating the derivative in \eqref{eq:block-stationarity} and
solving the resulting linear equation gives
\begin{equation}\label{eq:block-proportion-formula}
  \alpha_h:=
  \frac{\mathcal M_h(t_h^2)-\mathcal M_h(s_ht_h)}
  {\left\{\mathcal M_h(t_h^2)-\mathcal M_h(s_ht_h)\right\}+
  \left\{\mathcal M_h(s_h^2)-\mathcal M_h(s_ht_h)\right\}}.
\end{equation}
To analyze this quotient as $h\to0$, we factor $\mathcal M_h'$.  By
\eqref{eq:three-value-kkt}, the three values
$s_h^2,s_ht_h,t_h^2$ are zeros of $\mathcal M_h'$, and
\begin{equation}\label{eq:mh-derivative-factorization}
  \mathcal M_h'(z)
  =(z-s_h^2)(z-s_ht_h)(z-t_h^2)R(h,z),
\end{equation}
where $R$ is jointly analytic near $(0,r_*)$.
At $h=0$, the three zeros coincide at $r_*$, and direct calculation gives
\[
  \mathcal M_0'(z)
  =k_3(z-r_*)^3+O_d((z-r_*)^4),
  \qquad
  k_3:=\frac{d^5}{6(d-1)^2}>0.
\]
Thus $R(0,r_*)=k_3$.  Since every $z\in[s_h^2,t_h^2]$ satisfies
$|z-r_*|=O_d(|h|)$, joint analyticity of $R$ gives
\[
  R(h,z)=k_3+O_d(|z-r_*|+|h|)=k_3+O_d(|h|).
\]
Substituting this into \eqref{eq:mh-derivative-factorization}
and integrating over $[s_h^2,s_ht_h]$ and $[s_ht_h,t_h^2]$,
and using the expansions of $s_h$ and $t_h$ in powers of $h$,
show that the two differences in \eqref{eq:block-proportion-formula} vanish to order exactly four
and have the same leading term.  Canceling their common factor $h^4$ extends
the quotient analytically to $h=0$, where it equals $1/2$.  Finally,
$h\mapsto-h$ exchanges the two differences, so
$\alpha_{-h}=1-\alpha_h$.  The resulting analytic function $\alpha_h$ then
determines the analytic functions $q_h,r_h,\mu_h$, and the rank-one
factors $f_h$ and the bipodal graphons $W_h$.

The definitions of $q_h,r_h,\mu_h$ now give
$t(H,W_h)=q_h^v=r_h^m$.  The three equations
\eqref{eq:three-value-kkt} verify the stationarity condition
\eqref{eq:graphon-stationarity}, while
\eqref{eq:block-proportion-formula} verifies
\eqref{eq:block-stationarity}.
After decreasing $h_0$, continuity gives $0<p_h<r_h<1$,
$s_h,t_h,\alpha_h\in(0,1)$, and $\gamma_h,q_h>0$, hence $\mu_h>0$.
Thus, $(p_h,r_h)$ is admissible and $W_h$ is an interior bipodal graphon.

Finally, \Cref{lem:stationary-rank-one-bipodality} reduces any nearby rank-one graphon family in the
converse assertion to the same bipodal form that we selected in the beginning of the proof.
Repeating the construction, uniqueness in the implicit function theorem and in the linear
block-proportion equation forces the graphon, up to relabeling, to be
$W_h$ for a unique $h>0$, with multiplier $\mu_h$.
Full details are given in
\Cref{app:rank-one-kkt-family}.
\end{proof}

\begin{remark}[Dependence on $d$ and $H$]\label{rmk:rank-one-family-universality}
The equations \eqref{eq:three-value-kkt} and \eqref{eq:block-proportion-formula} depend only on $d$.
Thus, for fixed $d$, the local parameter curve $(p_h,r_h)$ and graphon
family $W_h$, up to relabeling, are independent of the particular
$d$-regular graph $H$.  The multiplier $\mu_h=\gamma_h/(m q_h^{v-2})$
and the interval on which we will establish global optimality and uniqueness
in the subsequent sections may depend on $H$.
\end{remark}
